% GENERATED FILE. Edit canonical lessons, then run npm run generate:books.
% canonical-source-hash: fnv1a64:e40ecd855fe8f2f7
% canonical-lessons: GU-C434-jhaptu, GU-C434-jharmar, GU-C434-reinkot, GU-C434-gheravu, GU-C434-bhijavu

\chapter{Showers and Drizzle}
\label{ch:gu-showers-and-drizzle}
\hlchaptermodality{434}

\begin{chapteropening}
\textbf{By the end of this chapter:} \emph{I can talk about rain coming and getting wet.}

Complete the last lesson of chapter 434: I can talk about rain coming and getting wet.
\end{chapteropening}

\section[\texorpdfstring{\gu{ઝાપટું} (jhāpṭũ)}{ઝાપટું (jhāpṭũ)}]{\gu{ઝાપટું} (jhāpṭũ) — a sharp shower (of rain)}

\label{lesson:GU-C434-jhaptu}



\begin{quote}
Before the new one: say the Gujarati for to suit, then the Gujarati for to shrink.
\end{quote}

\subsection*{You'll want to know: \gu{ઝાપટું}}
\textbf{\gu{ઝાપટું}} — \emph{jhāpṭũ} — \textquotedblleft{}a sharp shower (of rain)\textquotedblright{}.

\begin{quote}
\gu{બપોરે} \gu{વરસાદનું} \gu{એક} \gu{ઝાપટું} \gu{આવ્યું}.
\end{quote}

\begin{quote}
\emph{bapore varsādnũ ek jhāpṭũ āvyũ.}
\end{quote}

\begin{quote}
In the afternoon there was a sharp shower.
\end{quote}

A neuter noun.

\begin{cousinweb}
\textbf{\gu{ઝાપટું}} (\emph{jhāpṭũ}) goes back to \emph{*jhapp}, \textquotedblleft{}a sudden movement\textquotedblright{}, a form CDIAL reconstructs; CDIAL lists the Gujarati word, \textquotedblleft{}a sharp shower of rain\textquotedblright{} (CDIAL 5336).
\end{cousinweb}

\begin{grammarlens}[title={What you've built}]
One more word for getting around.
\end{grammarlens}

\subsection*{Your turn}
\begin{itemize}
\raggedright
  \item \emph{Say it:} \emph{jhāpṭũ}
  \item \emph{Say it:} \emph{jhāpṭũ}, once more
  \item \emph{Say it:} \emph{bapore varsādnũ ek jhāpṭũ āvyũ.} — In the afternoon there was a sharp shower.
  \item \emph{Recall:} say the Gujarati for to suit, then the Gujarati for to shrink, then say \emph{jhāpṭũ} again
\end{itemize}

\begin{tcolorbox}[breakable,colback=teal!4,colframe=teal!35!black,title={Before you move on}]
What does \gu{ઝાપટું} mean? (\textquotedblleft{}A sharp shower (of rain)\textquotedblright{}.) Say it once more.
\end{tcolorbox}
\section[\texorpdfstring{\gu{ઝરમર} (jharmar)}{ઝરમર (jharmar)}]{\gu{ઝરમર} (jharmar) — a drizzle}

\label{lesson:GU-C434-jharmar}



\begin{quote}
Before the new one: say the Gujarati for to shrink, then the Gujarati for a sharp shower.
\end{quote}

\subsection*{You'll want to know: \gu{ઝરમર}}
\textbf{\gu{ઝરમર}} — \emph{jharmar} — \textquotedblleft{}a drizzle\textquotedblright{}.

\begin{quote}
\gu{સવારથી} \gu{ઝરમર} \gu{વરસે} \gu{છે}.
\end{quote}

\begin{quote}
\emph{savārthī jharmar varse che.}
\end{quote}

\begin{quote}
It has been drizzling since the morning.
\end{quote}

A feminine noun.

\begin{cousinweb}
\textbf{\gu{ઝરમર}} (\emph{jharmar}) goes back to Sanskrit \emph{jhara}, \textquotedblleft{}a waterfall, a stream\textquotedblright{}; CDIAL lists the Gujarati word, \textquotedblleft{}a drizzle\textquotedblright{} (CDIAL 5343). Its doubled sound is like the rain itself.
\end{cousinweb}

\begin{grammarlens}[title={What you've built}]
One more word for getting around.
\end{grammarlens}

\subsection*{Your turn}
\begin{itemize}
\raggedright
  \item \emph{Say it:} \emph{jharmar}
  \item \emph{Say it:} \emph{jharmar}, once more
  \item \emph{Say it:} \emph{savārthī jharmar varse che.} — It has been drizzling since the morning.
  \item \emph{Recall:} say the Gujarati for to shrink, then the Gujarati for a sharp shower, then say \emph{jharmar} again
\end{itemize}

\begin{tcolorbox}[breakable,colback=teal!4,colframe=teal!35!black,title={Before you move on}]
What does \gu{ઝરમર} mean? (\textquotedblleft{}A drizzle\textquotedblright{}.) Say it once more.
\end{tcolorbox}
\section[\texorpdfstring{\gu{રેઈનકોટ} (reīnkoṭ)}{રેઈનકોટ (reīnkoṭ)}]{\gu{રેઈનકોટ} (reīnkoṭ) — a raincoat}

\label{lesson:GU-C434-reinkot}



\begin{quote}
Before the new one: say the Gujarati for a sharp shower, then the Gujarati for a drizzle.
\end{quote}

\subsection*{You'll want to know: \gu{રેઈનકોટ}}
\textbf{\gu{રેઈનકોટ}} — \emph{reīnkoṭ} — \textquotedblleft{}a raincoat\textquotedblright{}.

\begin{quote}
\gu{ચોમાસામાં} \gu{હું} \gu{રેઈનકોટ} \gu{પહેરીને} \gu{નીકળું} \gu{છું}.
\end{quote}

\begin{quote}
\emph{comāsāmā̃ hũ reīnkoṭ paherīne nīkḷũ chũ.}
\end{quote}

\begin{quote}
In the monsoon I put on a raincoat to go out.
\end{quote}

A masculine noun.

\begin{cousinweb}
\textbf{\gu{રેઈનકોટ}} (\emph{reīnkoṭ}) is borrowed from English \emph{raincoat}.
\end{cousinweb}

\begin{grammarlens}[title={What you've built}]
One more word for getting around.
\end{grammarlens}

\subsection*{Your turn}
\begin{itemize}
\raggedright
  \item \emph{Say it:} \emph{reīnkoṭ}
  \item \emph{Say it:} \emph{reīnkoṭ}, once more
  \item \emph{Say it:} \emph{comāsāmā̃ hũ reīnkoṭ paherīne nīkḷũ chũ.} — In the monsoon I put on a raincoat to go out.
  \item \emph{Recall:} say the Gujarati for a sharp shower, then the Gujarati for a drizzle, then say \emph{reīnkoṭ} again
\end{itemize}

\begin{tcolorbox}[breakable,colback=teal!4,colframe=teal!35!black,title={Before you move on}]
What does \gu{રેઈનકોટ} mean? (\textquotedblleft{}A raincoat\textquotedblright{}.) Say it once more.
\end{tcolorbox}
\section[\texorpdfstring{\gu{ઘેરાવું} (gherāvũ)}{ઘેરાવું (gherāvũ)}]{\gu{ઘેરાવું} (gherāvũ) — to gather (of clouds); to close in}

\label{lesson:GU-C434-gheravu}



\begin{quote}
Before the new one: say the Gujarati for a drizzle, then the Gujarati for a raincoat.
\end{quote}

\subsection*{You'll want to know: \gu{ઘેરાવું}}
\textbf{\gu{ઘેરાવું}} — \emph{gherāvũ} — \textquotedblleft{}to gather (of clouds); to close in\textquotedblright{}.

\begin{quote}
\gu{સાંજે} \gu{આકાશમાં} \gu{કાળાં} \gu{વાદળ} \gu{ઘેરાયાં}.
\end{quote}

\begin{quote}
\emph{sā̃je ākāśmā̃ kāḷā̃ vādaḷ gherāyā̃.}
\end{quote}

\begin{quote}
In the evening dark clouds gathered in the sky.
\end{quote}

Past (it …): \gu{તે} \gu{વાદળ} \gu{ઘેરાયું}, \emph{te vādaḷ gherāyũ}.

\begin{cousinweb}
\textbf{\gu{ઘેરાવું}} (\emph{gherāvũ}) goes back to \emph{*ghir}, \textquotedblleft{}to surround\textquotedblright{}, a form CDIAL reconstructs (CDIAL 4474): clouds that close in round the sky.
\end{cousinweb}

\begin{grammarlens}[title={What you've built}]
One more word for getting around.
\end{grammarlens}

\subsection*{Your turn}
\begin{itemize}
\raggedright
  \item \emph{Say it:} \emph{gherāvũ}
  \item \emph{Say it:} \emph{gherāvũ}, once more
  \item \emph{Say it:} \emph{sā̃je ākāśmā̃ kāḷā̃ vādaḷ gherāyā̃.} — In the evening dark clouds gathered in the sky.
  \item \emph{Recall:} say the Gujarati for a drizzle, then the Gujarati for a raincoat, then say \emph{gherāvũ} again
\end{itemize}

\begin{tcolorbox}[breakable,colback=teal!4,colframe=teal!35!black,title={Before you move on}]
What does \gu{ઘેરાવું} mean? (\textquotedblleft{}To gather (of clouds); to close in\textquotedblright{}.) Say it once more.
\end{tcolorbox}
\section[\texorpdfstring{\gu{ભીંજાવું} (bhī̃jāvũ)}{ભીંજાવું (bhī̃jāvũ)}]{\gu{ભીંજાવું} (bhī̃jāvũ) — to get wet through (in rain)}

\label{lesson:GU-C434-bhijavu}



\begin{quote}
Before the new one: say the Gujarati for a raincoat, then the Gujarati for to gather (of clouds).
\end{quote}

\subsection*{You'll want to know: \gu{ભીંજાવું}}
\textbf{\gu{ભીંજાવું}} — \emph{bhī̃jāvũ} — \textquotedblleft{}to get wet through (in rain)\textquotedblright{}.

\begin{quote}
\gu{છત્રી} \gu{ન} \gu{હતી}, \gu{એટલે} \gu{હું} \gu{ભીંજાઈ} \gu{ગયો}.
\end{quote}

\begin{quote}
\emph{chatrī na hatī, eṭle hũ bhī̃jāī gayo.}
\end{quote}

\begin{quote}
I had no umbrella, so I got wet through.
\end{quote}

Past (he …): \gu{તે} \gu{ભીંજાઈ} \gu{ગયો}, \emph{te bhī̃jāī gayo}.

\begin{cousinweb}
\textbf{\gu{ભીંજાવું}} (\emph{bhī̃jāvũ}) goes back to \emph{*bhiyañjati}, \textquotedblleft{}anoints, wets\textquotedblright{}, a form CDIAL reconstructs; CDIAL lists the Gujarati verb, \textquotedblleft{}to be wetted\textquotedblright{} (CDIAL 9503).
\end{cousinweb}

\begin{grammarlens}[title={What you've built}]
One more word for getting around.
\end{grammarlens}

\subsection*{Your turn}
\begin{itemize}
\raggedright
  \item \emph{Say it:} \emph{bhī̃jāvũ}
  \item \emph{Say it:} \emph{bhī̃jāvũ}, once more
  \item \emph{Say it:} \emph{chatrī na hatī, eṭle hũ bhī̃jāī gayo.} — I had no umbrella, so I got wet through.
  \item \emph{Recall:} say the Gujarati for a raincoat, then the Gujarati for to gather (of clouds), then say \emph{bhī̃jāvũ} again
\end{itemize}

\begin{tcolorbox}[breakable,colback=teal!4,colframe=teal!35!black,title={Before you move on}]
What does \gu{ભીંજાવું} mean? (\textquotedblleft{}To get wet through (in rain)\textquotedblright{}.) Say it once more.
\end{tcolorbox}
